\begin{lemma}[Value map on the deciding front]
For every locally constant $h$, there is a unique-value map
$h^{\dcl}:F^h\to E$ such that
\[
h^{\dcl}(s)=h(Y)\quad\text{whenever }s\segs Y.
\]
\end{lemma}
\begin{proof}
Existence of an enumeration extending any finite initial segment is
\noderef{FinitePrefixExtension}.  If two such enumerations extend the same $s$, the defining
conjunction in \noderef{DecidingFrontSet} says that $s$ is a deciding prefix.
By \noderef{DecidingPrefix}, any two enumerations having $s$ as a proper
initial segment therefore have equal $h$-values.  Hence choose any extension
and define $h^{\dcl}(s)$ by its value.  If $Y$ is another extension of $s$,
the same \noderef{DecidingPrefix} equality shows that this value is $h(Y)$.
Thus the definition is independent of the choice and has the displayed
property.
\end{proof}
